\begin{afterword}
\summaryhead

Some readers of ``Lotteries: A Waste of Hope'' objected that the author had no business criticizing other people's choices. The post answers that it is right for a human being to care about humanity's future, which includes the pursuit of truth; so other people's rationality is the author's business.

The idea is dangerous, the post grants: people have been burned for thinking wrongly. The remedy is to deny the ``therefore.'' One may judge that Susie said something wrong and argue against it, but must not set her on fire or silence her by violence or regulation. Science settles factual disputes by experiment rather than force, and the author wants the same fair fight for the whole future.

Relativists and advocates of selfishness, the post says, are not truly relativist or selfish, since they argue for their side; any true ones would be silent. Private self-deception often does little harm, but the author cannot help caring how other people think.

\responsehead

The rule at the center of the post is sound, and it is old. Argue against a belief, but do not compel the believer: Mill drew that line in \textsc{On Liberty} in 1859. The post presents it as ``my proposal.'' It answers the worry that caring about other people's beliefs leads to coercion, and the post states it plainly.

The answer to the opening question is a statement of the author's values. ``I believe that it is right and proper'' to care about the future; the pursuit of truth is part of that future; so others' rationality is the author's business. This explains why the author cares. It is not argued for. It is also the move that ``Why Truth?'' had described the previous November, where a sense of duty about truth leads one to think that others should look too, and where the author said ``I tend to be suspicious of morality as a motivation for rationality.'' The trouble named there was in one's own thinking. The safeguard offered here concerns only how others are treated.

The opponents are given in weaker forms than the positions they stand for. The view that no one should ever judge is given in exaggerated form and held by no one named. The comments that prompted the post question the author's standard or object to one judgment, which is a different objection. Relativists are said to be people who ``would not judge,'' though the main relativist positions in philosophy say that moral truth is relative, not that judging is forbidden. The charge against advocates of selfishness is a known objection to ethical egoism, which shows that such advocacy is unusual, not that it is inconsistent. The paragraph after makes the claim safe from counterexample: any relativist or egoist who speaks has already been counted as not a true one.

The ending is candid about scope. For private belief in what is pleasant, the post's case reduces to ``a small darkness only'' and the author's mourning.

In short: an old liberal rule, argue but do not compel, presented as the author's proposal, attached to a statement of the author's values, and defended against relativists and egoists defined so that no true one can be heard.
\end{afterword}
